\chapter{Portability Context and Research Roadmap}
\label{app:portability-roadmap}

The battery thesis is intentionally not a universal proof document.  The
repo nevertheless contains enough adapter structure and roadmap material to
justify one bounded appendix on how the battery result interfaces with a
larger ORIUS program.  This appendix serves that purpose.

\section{From Battery Adapter to a Portable Runtime Pattern}

The battery result is structurally interesting because the DC3S pipeline
operates on a small set of abstract objects:

\begin{enumerate}
  \item an observation vector,
  \item a reliability score,
  \item an uncertainty region,
  \item a feasible action set, and
  \item a repair or fallback operator.
\end{enumerate}

The battery adapter is one concrete instantiation of those objects.  The
abstract pattern is portable; the evidence depth is not.  That distinction is
exactly why the battery thesis can mention adapters without turning into a
universal guarantee document.

\section{Adapter Mapping}

\begin{table}[ht]
\centering
\caption{Battery objects and their domain-general counterparts.}
\label{tab:portable-adapter-map}
\begin{tabular}{lll}
\toprule
\textbf{ORIUS object} & \textbf{Battery instantiation} & \textbf{Portable meaning} \\
\midrule
$z_t$ & SCADA packet (SOC, voltage, current, load) & raw telemetry vector \\
$w_t$ & OQE composite score & observation reliability \\
$C_t^{\mathrm{RAC}}$ & SOC conformal interval & state uncertainty region \\
$\mathcal{A}_t$ & charge/discharge bounds & feasible action set \\
$a_t^{\mathrm{safe}}$ & SOC-preserving dispatch & safety-projected action \\
\bottomrule
\end{tabular}
\end{table}

\section{What Changes by Domain}

Porting ORIUS to a new domain requires replacing four domain-specific
components:

\begin{description}
  \item[Plant dynamics.] Battery uses an SOC integrator; other domains may be
    nonlinear or high-dimensional.
  \item[State definition.] SOC is scalar; vehicles, industrial systems, and
    aerospace plants carry richer state.
  \item[Constraint geometry.] Batteries are dominated by box-style limits;
    other domains may need polytopic or nonlinear repair.
  \item[Fault taxonomy.] Battery dropout, drift, stale sensors, and spikes
    must be re-instantiated for the target telemetry stack.
\end{description}

\section{What the Repo Already Contains}

\begin{table}[ht]
\centering
\caption{Current bounded portability surfaces present in the repo.}
\label{tab:portability-surfaces}
\begin{tabular}{p{2.5cm}p{3.2cm}p{5.1cm}}
\toprule
\textbf{Surface} & \textbf{Primary path} & \textbf{Bounded role} \\
\midrule
Battery adapter & \path{src/orius/dc3s/battery_adapter.py} & Locked proof and evidence domain \\
Vehicle adapter & \path{src/orius/vehicles/vehicle_adapter.py} & Prototype second-domain portability surface \\
ORIUS-Bench tracks & \path{src/orius/orius_bench/} & Shared evaluation contract \\
Universal registry & \path{src/orius/universal_framework/domain_registry.py} & Adapter wiring and dispatch \\
CertOS runtime & \path{src/orius/certos/} & Runtime pattern with bounded portability evidence \\
\bottomrule
\end{tabular}
\end{table}

\section{Leaderboards as Portability Evidence, Not Dominance Evidence}

The broader ORIUS-Bench leaderboard provides context for portability.  It
does \emph{not} carry the same proof weight as the battery domain.

\begin{table}[ht]
\centering
\caption{Vehicle-domain excerpt from the shared ORIUS-Bench leaderboard.}
\label{tab:vehicle-leaderboard-excerpt}
\begin{tabular}{lrrrr}
\toprule
\textbf{Controller} & \textbf{Seed} & \textbf{TSVR} & \textbf{GDQ} & \textbf{IR} \\
\midrule
nominal & 1000 & 0.0833 & 0.9167 & 0.4583 \\
nominal & 1001 & 0.0833 & 0.9167 & 0.3333 \\
robust  & 1000 & 0.0000 & 1.0000 & 0.4583 \\
robust  & 1001 & 0.0000 & 1.0000 & 0.3333 \\
dc3s    & 1000 & 0.1667 & 0.8333 & 0.4583 \\
dc3s    & 1001 & 0.1667 & 0.8333 & 0.3333 \\
fallback& 1000 & 0.0000 & 1.0000 & 0.4583 \\
fallback& 1001 & 0.0000 & 1.0000 & 0.3333 \\
\bottomrule
\end{tabular}
\end{table}

This excerpt is useful because it shows the shared contract actually runs on
another domain.  It is not useful as evidence of universal DC3S dominance,
and the thesis should not read it that way.

\section{Research Roadmap}

The long-term program remains dependency-ordered:

\begin{enumerate}
  \item preserve and defend the locked battery foundation,
  \item harden remaining battery gaps before broadening claims,
  \item extract domain-independent ORIUS core modules cleanly,
  \item validate one additional domain under the full dual-trajectory
    protocol,
  \item strengthen hardware-in-the-loop and controller-in-the-loop evidence,
  \item release the benchmark and adapter contract more publicly, and
  \item only then widen runtime and deployment claims.
\end{enumerate}

\section{Deployment Sequence}

The practical deployment ladder for the battery stack is:

\begin{enumerate}
  \item benchmark replay with observed-vs.-true trajectory separation,
  \item controller-in-the-loop runtime rehearsal with complete
    certificates,
  \item hardware-in-the-loop with real telemetry transport and relay
    behavior,
  \item supervised pilot operation under bounded fallback policies, and
  \item wider deployment only after those layers remain manifest-traceable.
\end{enumerate}

\section{Why This Appendix Exists}

This appendix exists to make one distinction explicit.  The thesis is long
because it contains a complete battery proof surface plus a disciplined
program of bounded extensions.  It is \emph{not} long because it claims to
have solved every future domain already.
